\documentclass[12pt]{article}
\usepackage[utf8]{inputenc}
\usepackage[T1]{fontenc}
\usepackage{hyperref}
\title{L\"oy\"op Language Resources for K-12 Curriculum Development}
\author{Compiled by Claude Sonnet 5 (claude-sonnet-5) for Andr\'as Kornai}
\date{}
\begin{document}
\maketitle

\section{Introduction}
L\"oy\"op (ISO 639-3 urr) is a small Oceanic language spoken on the east coast of Ureparapara in the Banks Islands, Vanuatu. Estimates range from about 200 to 250 speakers. It was formerly called Lehalurup, and it is distinct from the neighbouring Lehali language on the west coast of the island.

\section{Primary Educational Resources}
No primers, textbooks or school programmes were found in search results.

\section{Linguistic and Reference Materials}
\begin{itemize}
\item Wikipedia entry: \url{https://en.wikipedia.org/wiki/L%C3%B6y%C3%B6p_language}
\item Omniglot page on L\"oy\"op, which notes Codrington's 1885 documentation: \url{https://www.omniglot.com/writing/loyop.htm}
\item Neighbouring language Lehali, for comparison: \url{https://en.wikipedia.org/wiki/Lehali_language}
\end{itemize}
No dictionary or grammar was found.

\section{Religious and Community Resources}
None found.

\section{Digital and Audio Resources}
The Endangered Languages Project has a listing, and Verbix has an endangered-languages page: \url{https://www.endangeredlanguages.com/lang/3575} and \url{https://docs.verbix.com/EndangeredLanguages/Loyop}

\section{Curriculum Development Context}
Only overview material exists. A curriculum would need new documentation work with the community and Vanuatu education authorities.
\end{document}